
\documentclass[11pt]{article}
\usepackage[a4paper,margin=0.95in]{geometry}
\usepackage{fontspec}
\setmainfont{Noto Serif}
\usepackage{microtype}
\usepackage{amsmath,amssymb,mathtools}
\usepackage{csquotes}
\usepackage{array}
\usepackage{hyperref}
\hypersetup{colorlinks=false,pdfborder={0 0 0}}
\setlength{\parindent}{1.15em}
\setlength{\parskip}{0.18em}
\allowdisplaybreaks
\sloppy
\newcommand{\scn}[1]{\textsc{#1}}
\newcommand{\arctg}{\operatorname{arctg}}
\begin{document}
\begin{center}
{\Large\bfseries NOTE SUR UN THÉORÈME D'ANALYSE INDÉTERMINÉE\par}
\vspace{0.8em}
\end{center}

Je lis dans le résumé des séances de l'Académie des Sciences, inséré dans les Annales de Chimie et de Physique: «M. Lamé présente un mémoire sur l'impossibilité de l'équation $x^5+y^5=2^n z^5$, en nombres entiers». Cela me rappelle un théorème du même genre sur lequel je suis tombé, il y a plus d'un an, et dont voici l'énoncé.

L'équation $x^3\pm y^3=2^n z^3$ est impossible en nombres entiers et positifs.

Pour démontrer ce théorème, il suffira de prouver l'impossibilité de ces trois équations
\[
\tag{a}
 x^3\pm y^3=z^3,\qquad x^3\pm y^3=2z^3,\qquad x^3\pm y^3=4z^3.
\]
Car, si l'on voulait supposer $x^3\pm y^3=2^n z^3$, $n$ étant $=0$ ou $>3$, on n'aurait qu'à faire $n=3k+r$, $r$ étant un des nombres $0,1,2$; ce qui changerait l'équation précédente en celle-ci $x^3\pm y^3=2^r(2^kz)^3$, qui rentre dans une des équations (a) supposées impossibles.

L'impossibilité des deux premières des équations (a) est démontrée depuis longtemps, comme on peut le voir dans le second volume de l'Algèbre d'\scn{Euler} ou dans la théorie des nombres de M. \scn{Legendre}. Je n'ai donc qu'à m'occuper de la troisième.

Soit, s'il est possible, $x^3+y^3=4z^3$. Si, comme il est permis, on suppose les nombres $x,y$ premiers entre eux, ils seront impairs tous deux. Faisant donc $x=p+q$ et $\pm y=p-q$, les nombres $p$ et $q$ seront entiers, premiers entre eux, et de plus l'un pair, l'autre impair. La substitution des valeurs précédentes donne l'équation
\[
 p(p^2+3q^2)=2z^3,
\]
qui montre que c'est $p$ qu'il faut supposer pair, sans quoi le premier membre serait impair; $\frac p2$ est donc entier et l'on a
\[
 \frac p2(p^2+3q^2)=z^3.
\]
Il y a maintenant deux cas à distinguer selon que $p$ est ou n'est pas divisible par $3$.

Premier cas. -- Si $p$ n'est pas divisible par $3$, $\frac p2$ et $p^2+3q^2$ seront premiers entre eux, et leur produit étant un cube, chacun de ces nombres devra en être un. Pour faire $p^2+3q^2$ égal à un cube de la manière la plus générale, on n'a, comme on sait, qu'à poser ces deux équations
\[
 p=m^3-9mn^2,\qquad q=3m^2n-3n^3.
\]
D'après ce qui a été dit sur les nombres $p$ et $q$, on verra facilement que $m$ doit être premier à $n$, et de plus pair et non divisible par $3$. Il suit de là que les trois facteurs du second membre de l'équation
\[
 \frac p2=\frac m2(m+3n)(m-3n)
\]
sont premiers entre eux. Leur produit $\frac p2$ devant être un cube, il faudra que chacun de ces facteurs en soit un; on aura par conséquent
\[
 \frac m2=\alpha^3,\qquad m+3n=\beta^3,\qquad m-3n=\gamma^3,
\]
d'où l'on tire cette équation $\beta^3+\gamma^3=4\alpha^3$, qui est semblable à la première $x^3+y^3=4z^3$, mais exprimée en nombres beaucoup plus petits.

Second cas. Si $p$ est divisible par $3$, on fera $p=3r$, et $r$ sera, comme $p$, pair et premier à $q$. On aura ainsi
\[
 \frac{9r}{2}(q^2+3r^2)=z^3.
\]
On verra facilement que les deux facteurs du premier membre sont premiers entre eux; il faut donc que chacun d'eux soit un cube. Pour faire $q^2+3r^2$ égal à un cube, on posera, comme dans le premier cas,
\[
 q=m^3-9mn^2,\qquad r=3m^2n-3n^3.
\]
On peut s'assurer que $n$ doit être premier à $m$, et de plus pair; $\frac{9r}{2}$ sera donc entier et l'on aura
\[
 \frac{9r}{2}=27\frac n2(m+n)(m-n).
\]
Les facteurs $\frac n2$, $m+n$, $m-n$ sont premiers entre eux d'après ce qu'on vient de dire sur $m$ et $n$; et leur produit multiplié par le cube $27$ devant être un cube, il faut que chacun d'eux soit un cube. On aura donc
\[
 \frac n2=\alpha^3,\qquad m+n=\beta^3,\qquad m-n=\gamma^3,
\]
et par conséquent $\beta^3-\gamma^3=4\alpha^3$, équation semblable à la proposée $x^3\pm y^3=4z^3$, mais exprimée en nombres beaucoup plus petits.

On conclut de là, à la manière ordinaire, que l'équation proposée n'est pas possible.

\end{document}
